\documentclass[10pt]{article}

\usepackage[margin=1in]{geometry}
\usepackage{setspace}
\usepackage{parskip}
\usepackage{hyperref}
\usepackage{titling}
\usepackage{amsmath, amssymb}
\usepackage[numbers]{natbib}

\title{\textbf{Prompt Degrees of Freedom in LLM-Based Social Science: A Specification Curve Framework for Transparent Reporting}}

\author{%
\small Jacob Crainic$^{1}$\\[4pt]
{\small $^{1}$~Management Sciences Lab, Yee Collins Research Group}\\[2pt]
{\small \texttt{jacobcrainic2008@gmail.com}}%
}

\date{}

\begin{document}
\maketitle

\begin{abstract}
When researchers use large language models as synthetic survey respondents, the prompt is an analytic choice. Model selection, persona wording, question framing, system instructions, temperature, and few-shot exemplars each represent decision points that may shift substantive conclusions, yet current practice reports results from a single configuration and treats the rest as incidental. This is the prompt-domain version of the ``garden of forking paths'' that has driven the credibility revolution in psychology and economics \cite{gelman2014statistical, simonsohn2020specification}. We introduce \textit{Prompt Specification Curve Analysis} (P-SCA): a method that maps how conclusions vary across the space of defensible prompt designs, decomposes outcome variance by prompt dimension, and produces a reporting standard for LLM-based social science.

Argyle et al.\ \cite{argyle2023out} showed that demographically conditioned language models can approximate U.S.\ public opinion distributions, a paradigm they call ``silicon sampling.'' But fidelity has limits. Bisbee et al.\ \cite{bisbee2024synthetic} found that it degrades for underrepresented populations. Sclar et al.\ \cite{sclar2024quantifying} documented that minor prompt perturbations, including whitespace and label ordering, can shift NLP benchmark accuracy by up to 76 points. Karell and Davidson \cite{karell2025integrating} identify prompt design as one of three core methodological challenges for generative AI in social science, and note that the field has no established approach to treating prompting as an experimental condition. Prompt sensitivity is real, but the literature on it remains descriptive. There is no method that lets researchers enumerate defensible prompt designs for a given study, measure how much each design dimension contributes to outcome variation, or test whether a reported finding holds across the prompt multiverse or depends on a particular configuration. Specification curve analysis \cite{simonsohn2020specification} solved this problem for traditional analytic choices; Steegen et al.\ \cite{steegen2016increasing} generalized it as ``multiverse analysis,'' in which every defensible analytic path is run and the distribution of outcomes reported. Our \textit{prompt multiverse} borrows this architecture and applies it to the prompt design space.

We define a six-dimensional \textit{prompt multiverse} over the design choices in LLM-based opinion elicitation, spanning model architecture (GPT-5.4, GPT-5.4-nano, Claude Sonnet 4.6, Gemini 2.5 Flash), persona format (bare demographic tag, first-person narrative, third-person vignette, survey-style respondent block), question framing (direct elicitation, Likert scale, forced choice), system prompt register (neutral, academic researcher, survey administrator), sampling temperature (0.0, 0.3, 0.7, 1.0), and few-shot exemplar count (0, 1, 3). Researchers make these choices routinely but rarely justify or report them. We sample 600 specifications from this $4 \times 4 \times 3 \times 3 \times 4 \times 3 = 1{,}728$-cell space using Latin Hypercube Sampling \cite{mckay1979comparison}, covering the full multiverse without exhaustive enumeration. Each specification is administered to 20 demographic profiles drawn from five battleground states (PA, MI, WI, AZ, GA), balanced by party identification, with 5 repeated draws per profile to estimate within-specification variance. For each specification, we compute the partisan opinion gap (mean Democratic response minus mean Republican response) on three American National Election Studies items: government spending, immigration levels, and gun regulation. If a finding holds in 95\% of specifications, it is likely a property of the underlying opinion distribution rather than an artifact of prompt construction. If it holds in 55\%, the result depends on design choices the researcher did not report.

To identify which dimensions drive variation, we compute $\eta^2$ from one-way decompositions of response scores by each dimension, producing a rank-ordering of design choices by their contribution to outcome variation. The full study extends this with Sobol sensitivity indices \cite{saltelli2008global} estimated via Saltelli's quasi-random scheme, which captures interaction effects between dimensions. Joint inference follows the permutation-based procedure of \cite{simonsohn2020specification}, testing whether LLM partisan structure exceeds what a null model of prompt-driven noise would produce. We plot specification curves sorted by effect magnitude, following the graphical conventions of the original specification curve literature.

This study provides the first specification curve analysis of prompt design in a social science application. In pilot runs on the government spending item (100 specifications, 4 models, 20 profiles), the expected partisan gap holds in 71\% of specifications. Model architecture dominates the variance decomposition ($\eta^2 = 0.34$), question framing is second ($\eta^2 = 0.18$), and persona format, system prompt, temperature, and few-shot count each contribute under 3\%. Model choice matters roughly 25 times more than persona wording. These preliminary results give researchers concrete guidance on where robustness checks are necessary and where they are not. We also propose a \textit{Prompt Specification Reporting Standard}: authors document their prompt multiverse, include specification curve plots, and disclose the share of defensible specifications consistent with their claimed finding. This responds to the transparency gap that Karell and Davidson \cite{karell2025integrating} identified, and gives reviewers a concrete criterion for evaluating LLM-based empirical work. Because P-SCA treats model version as an explicit dimension, it also tracks how conclusions shift across model generations. A study run on GPT-5.4 today and replicated on its successor next year can be compared within the same specification curve, making replication legible even as the underlying systems change. The method is not specific to political opinion; the same pipeline applies wherever LLMs serve as behavioral proxies. All prompt templates, sampling code, API call logs, and analysis scripts will be released as an open-source toolkit.
\end{abstract}

\begin{thebibliography}{10}
\small\setlength{\itemsep}{1pt}

\bibitem{argyle2023out}
Argyle, L. P., Busby, E. C., Fulda, N., Gubler, J. R., Rytting, C., and Wingate, D. (2023).
\textit{Out of One, Many: Using Language Models to Simulate Human Samples.}
Political Analysis, 31(3), 337--351.

\bibitem{bisbee2024synthetic}
Bisbee, J., Clinton, J. D., Dorff, C., Kenkel, B., and Larson, J. (2024).
\textit{Synthetic Replacements for Human Survey Data? The Perils of Large Language Models.}
Political Analysis, 32(4), 401--416.

\bibitem{gelman2014statistical}
Gelman, A. and Loken, E. (2014).
\textit{The Statistical Crisis in Science.}
American Scientist, 102(6), 460--465.

\bibitem{karell2025integrating}
Karell, D. and Davidson, T. (2025).
\textit{Integrating Generative Artificial Intelligence into Social Science Research: Measurement, Prompting, and Simulation.}
Sociological Methods \& Research.

\bibitem{mckay1979comparison}
McKay, M. D., Beckman, R. J., and Conover, W. J. (1979).
\textit{A Comparison of Three Methods for Selecting Values of Input Variables in the Analysis of Output from a Computer Code.}
Technometrics, 21(2), 239--245.

\bibitem{saltelli2008global}
Saltelli, A., Ratto, M., Andres, T., Campolongo, F., Cariboni, J., Gatelli, D., Saisana, M., and Tarantola, S. (2008).
\textit{Global Sensitivity Analysis: The Primer.}
John Wiley \& Sons.

\bibitem{sclar2024quantifying}
Sclar, M., Choi, Y., Tsvetkov, Y., and Suhr, A. (2024).
\textit{Quantifying Language Models' Sensitivity to Spurious Features in Prompt Design with a Focus on Social Science Tasks.}
Proceedings of ACL 2024.

\bibitem{simonsohn2020specification}
Simonsohn, U., Simmons, J. P., and Nelson, L. D. (2020).
\textit{Specification Curve Analysis.}
Nature Human Behaviour, 4(11), 1208--1214.

\bibitem{steegen2016increasing}
Steegen, S., Tuerlinckx, F., Gelman, A., and Vanpaemel, W. (2016).
\textit{Increasing Transparency Through a Multiverse Analysis.}
Perspectives on Psychological Science, 11(5), 702--712.

\end{thebibliography}

\end{document}
